\section{A scalar Fourier-coefficient lemma}\label{sec:coefficient}

For a function on a sign cube, write
\[
\Odd h(z)=\frac{h(z)-h(-z)}2,
\]
where \(-z\) negates every coordinate. We call \(h\) odd if \(\Odd h=h\).
Fourier coefficients always use uniform measure. For an \(m\times m\)
sign matrix \(S\) and a binary matrix \(P\in\{0,1\}^{m\times m}\), put
\[
S^P=\prod_{i,j}S_{ij}^{P_{ij}},\qquad
\widehat h(P)=\E_S h(S)S^P.
\]
For vector cubes we use binary indices \(p\in\{0,1\}^m\);
\(|p|\) denotes the number of ones and \(e_i\) the \(i\)-th unit vector.
The notation \(\sigma S\) multiplies row \(i\) by \(\sigma_i\), and
\(S\tau\) multiplies column \(j\) by \(\tau_j\).
We denote row \(i\) and column \(j\) by \(S_{i,:}\) and \(S_{:,j}\), respectively.
All ranks in this section are over \(\F_2\).

The point of the binary rank is the identity
\[
\prod_{i,j}\chi(u_i,v_j)^{P_{ij}}
=(-1)^{\sum_{i,j}P_{ij}u_i\cdot v_j},
\qquad u_i,v_j\in\F_2^n.
\]
An index of rank \(k\) yields \(k\) independent bilinear terms after
binary changes of variables. In Section~\ref{sec:walsh}, this will turn a
nonzero coefficient into a Walsh pairing in \(k\) times as many
coordinates. We therefore need a coefficient whose index has rank
at least two.

\par\medskip\noindent\begin{minipage}{\textwidth}
For example, take \(S_{ij}=\chi(u_i,v_j)\) for \(1\le i,j\le3\).
The two indices with an odd number of ones
\[
\begin{gathered}
P_1=\underbrace{\begin{pmatrix}
1&1&1\\1&1&1\\1&1&1
\end{pmatrix}}_{\text{binary rank }1},
\qquad S^{P_1}=\chi(u_1+u_2+u_3,v_1+v_2+v_3),\\[6pt]
P_2=\underbrace{\begin{pmatrix}
1&1&0\\0&1&0\\0&0&0
\end{pmatrix}}_{\text{binary rank }2},
\qquad S^{P_2}=\chi(u_1,v_1+v_2)\chi(u_2,v_2)
\end{gathered}
\]
yield one and two Walsh factors, respectively: \(v_1+v_2\) and
\(v_2\) vary independently. \end{minipage}\par\medskip
Rank-one indices have rectangular support.
The next lemma rules out an expansion supported only on odd rectangles
when both families of prescribed moments are large enough.

\begin{lemma}[A coefficient of rank at least two]\label{lem:coefficient}
Let \(M_0=4\), and suppose that
\[
4\le d=\lceil2\log_2m\rceil<m.
\]
Let \(a,b>0\), put \(A=ma\), \(B=mb\), and assume
\begin{equation}\label{eq:large-parameters}
\min(A,B)\ge20M_0d(d+1)^2.
\end{equation}
There is \(c_*>0\), depending only on \(m,a,b\), with the following
property. Suppose that
\[
\Psi:\{-1,1\}^{m\times m}\to\R,\qquad
f,g:\{-1,1\}^m\to\R
\]
are odd, have sup norms at most \(M_0\), and satisfy
\begin{equation}\label{eq:scalar-first}
\widehat f(e_i)=a,\qquad \widehat g(e_j)=b
\end{equation}
together with the projection identities
\begin{equation}\label{eq:scalar-projections}
\begin{split}
\E_\sigma\sigma_i\Psi(\sigma S)&=a\,g(S_{i,:}),\\
\E_\tau\tau_j\Psi(S\tau)&=b\,f(S_{:,j})
\end{split}
\end{equation}
for all \(i,j,S\). Then
\[
|\widehat\Psi(P)|\ge c_*
\]
for some binary matrix \(P\) of rank at least two.
\end{lemma}

\begin{proof}
We first reduce the uniform lower bound to a nonvanishing assertion.
If every index in the Fourier support of \(\Psi\) had rank at most one, a damped polynomial would encode
their rectangular supports. Its prescribed mixed derivative will
contradict a bound obtained by splitting it into two parts, each of
low degree in one variable.

The admissible triples \((\Psi,f,g)\) form a compact set in a
finite-dimensional space. The maximum of
\(|\widehat\Psi(P)|\) over \(\rank P\ge2\) is continuous.
It therefore suffices to show that this maximum never vanishes.
If the admissible set is empty, any positive \(c_*\) will do.

Suppose, to the contrary, that the Fourier support of \(\Psi\)
has rank at most one. Oddness eliminates indices with an even
number of ones. Every remaining index has a unique form
\(P=ph^t\) with \(p,h\ne0\), and both \(k=|p|\) and \(l=|h|\) are odd.
Let
\[
C_{kl}=\sum_{\substack{|p|=k\\|h|=l}}\widehat\Psi(ph^t).
\]
Parseval's identity and Cauchy--Schwarz give
\begin{equation}\label{eq:Ckl-bound}
|C_{kl}|\le M_0\sqrt{\binom mk\binom ml}.
\end{equation}
To suppress rectangles that are large in both directions, consider
the damped polynomial
\[
H(s,u)=\sum_{\substack{1\le k,l\le m\\ k,l\ {\rm odd}}}
C_{kl}s^ku^l2^{-kl}.
\]
It has sup norm at most \(M_0\) on \([-1,1]^2\).
To see this, take independent signs \(D_i,E_j,R_{ij}\) with
means \(s,u,1/2\), respectively, and set \(S_{ij}=D_iE_jR_{ij}\).
For an odd rectangle \(ph^t\),
\[
\E S^{ph^t}=s^{|p|}u^{|h|}2^{-|p||h|},
\]
so \(H(s,u)=\E\Psi(S)\).

Delete from \(H\) the terms with both \(k>d\) and \(l>d\), obtaining
\(H_0\). If \(k,l>d\), then
\[
(k+l)\log_2m\le kl,\qquad kl/2\ge k+l.
\]
Using \eqref{eq:Ckl-bound}, the deleted part is bounded in sup norm by
\[
M_0\sum_{k,l>d}m^{(k+l)/2}2^{-kl}
\le M_0\sum_{k,l>d}2^{-kl/2}
\le M_0\sum_{k,l>d}2^{-k-l}\le M_0.
\]
Consequently \(\|H_0\|_\infty\le2M_0\).

For an odd rectangle \(ph^t\), the vectors of row and column parities
are \(p\) and \(h\). Comparing Fourier coefficients in
\eqref{eq:scalar-projections} therefore gives
\[
C_{1l}=ma\sum_{|h|=l}\widehat g(h),\qquad
C_{k1}=mb\sum_{|p|=k}\widehat f(p).
\]
The deletion defining \(H_0\) does not affect terms with \(k=1\) or
\(l=1\). Hence
\begin{equation}\label{eq:H-derivatives}
\begin{split}
\partial_sH_0(0,u)&=A\sum_h\widehat g(h)(u/2)^{|h|},\\
\partial_uH_0(s,0)&=B\sum_p\widehat f(p)(s/2)^{|p|},\\
\partial_s\partial_uH_0(0,0)&=AB/2.
\end{split}
\end{equation}
Independent signs of common mean \(u/2\), or \(s/2\), show that
the first two expressions have sup norms at most \(AM_0\) and \(BM_0\).

Every retained monomial has low degree in at least one variable.
We separate these two possibilities with controlled sup norms.
Expand \(H_0\) in the Chebyshev basis \(T_j(s)\), defined by
\(T_j(\cos\theta)=\cos(j\theta)\). Let \(H_1\) retain the terms
with \(j\le d\), and set \(H_2=H_0-H_1\).
Then
\begin{equation}\label{eq:separate-degrees}
\deg_sH_1\le d,\qquad \deg_uH_2\le d.
\end{equation}
Indeed, every monomial of \(H_0\) has \(s\)-degree at most \(d\)
or \(u\)-degree at most \(d\), and truncation in the Chebyshev
basis annihilates every monomial of the former kind in \(H_2\).

Each Chebyshev coefficient of a polynomial on \([-1,1]\) is bounded
by twice its sup norm, by cosine orthogonality. Since \(|T_j|\le1\),
\[
\|H_1\|_\infty\le4M_0(d+1),\qquad
\|H_2\|_\infty\le6M_0(d+1).
\]
Also, \(|T_j'(0)|\le j\), so every polynomial \(Q\) of degree at most
\(d\) satisfies
\begin{equation}\label{eq:derivative-bound}
|Q'(0)|\le d(d+1)\|Q\|_\infty.
\end{equation}
First apply \eqref{eq:derivative-bound} in \(s\) to \(H_1\), and
in \(u\) to \(H_2\). Then \eqref{eq:H-derivatives} implies
\[
\begin{split}
\|\partial_sH_2(0,\cdot)\|_\infty
&\le AM_0+4M_0d(d+1)^2,\\
\|\partial_uH_1(\cdot,0)\|_\infty
&\le BM_0+6M_0d(d+1)^2.
\end{split}
\]
By \eqref{eq:separate-degrees}, these are polynomials of degree
at most \(d\) in their displayed variables. Applying
\eqref{eq:derivative-bound} once more and adding yields
\begin{equation}\label{eq:mixed-upper}
|\partial_s\partial_uH_0(0,0)|
\le M_0d(d+1)\bigl(A+B+10d(d+1)^2\bigr).
\end{equation}
On the other hand, \eqref{eq:large-parameters} gives
\[
\frac{AB}{2}\ge5M_0d(d+1)^2(A+B),\qquad
10d(d+1)^2\le\frac{A+B}{4M_0}.
\]
Since \(5(d+1)>1+1/(4M_0)\), these inequalities contradict
\eqref{eq:H-derivatives} and \eqref{eq:mixed-upper}.
\end{proof}
